\documentclass{article}
\usepackage[T1]{fontenc}
\usepackage{lmodern}
\usepackage{graphicx}
\usepackage{amsmath,amssymb}
\usepackage[a4paper,margin=2cm]{geometry}
\usepackage[absolute,overlay]{textpos}
\setlength{\TPHorizModule}{1mm}
\setlength{\TPVertModule}{1mm}
\usepackage[indonesian]{babel}
\usepackage{hyperref}
\setlength{\emergencystretch}{2em}
% unit: Halaman 24

\begin{document}

% === Halaman 24 (python/native) ===

24

• Hasil iterasi ke-2:

 hereTCSWPeYraWHarSseQithaYraOeaWHstatePFairaWHKrSTYhtm

etheKeetPeJrSmaYPassGaseiWQhiGhitQaseWGPSseHitQasaKeaT

tiJTPsGaraKaeTsaWHatthattimeTWAWSQWtSWatTraPistsSJGSTr

seaYreatCriUeiWasGieWtiJiGCSiWtSJOieQthereQeretQSrSTWH

KPaGAsCStsWearSWeeBtremitFSJtheKaGAaWHaPSWYSWeWeartheS

therthesGaPesQereeBGeeHiWYPFharHaWHYPSssFQithaPPtheaCC

earaWGeSJKTrWisheHYSPHtheQeiYhtSJtheiWseGtQasOerFremar

AaKPeaWHtaAiWYaPPthiWYsiWtSGSWsiHeratiSWiGSTPHharHPFKP

ameNTCiterJSrhisSCiWiSWresCeGtiWYit

• Teruskan, dengan menerka kata-kata yang sudah dikenal, misalnya 

remarA mungkin remark , dsb

\end{document}
